\documentclass[../numerical_methods.tex]{subfiles}
\begin{document}
\section{Aula 03 - 10 de Março, 2026}
\subsection{Motivações}
\begin{itemize}
	\item Vetores e listas;
\end{itemize}
\subsection{Vetores e listas}
Os objetivos da aula de hoje incluem entender o que são listas, conhecidos também como vetores, incluindo como acessar e modificar elementos, percorrer listas com loops, realizar operações básicas e resolver alguns desafios.

Suponhamos que nosso objetivo é armazenar as notas de vários alunos; seria muito trabalhoso associar uma variável a cada um deles, como em
\begin{minted}[bgcolor = DarkBackground, linenos=true]{python}
  nota1=7.5
  nota2=8.0
  nota3=6.5
  nota4=9.0
  nota5=7.0

  print(nota1, nota2, nota3, nota4, nota5 )
\end{minted}
Para esse tipo de problema, utilizamos uma lista, demarcada por \texttt{}, e cada item é um índice dela, começando em 0:
\begin{minted}[bgcolor = DarkBackground, linenos=true]{python}
  notas = [7.5, 8.0, 6.5, 9.8, 7.0]

  print(notas[0])
  print(notas[2])
\end{minted}

Podemos manipular diretamente os elementos apenas sinalizando que certo índice recebe um novo valor:
\begin{minted}[bgcolor = DarkBackground, linenos=true]{python}
  notas = [7.5, 8.0, 6.5, 9.8, 7.0]

  notas([2])=9.9

  print(notas[2])
\end{minted}

Obter o tamanho de uma lista qualquer pode ser feito usando o comando \texttt{len()}, e percorrê-las é tão simples quanto usar um for:
\begin{minted}[bgcolor = DarkBackground, linenos=true]{python}
  notas = [7.5, 8.0, 6.5, 9.8, 7.0]

  print(len(notas))

  for nota in notas:
      print(nota)

  # Outra forma de fazer usando o len:

  for indice in range(0, len(notas)):
      print(nota) 
\end{minted}

Juntanto com o que vimos até o momento, fica tranquilo entender como calcular a média da turma

\begin{listing}
	\caption{Média das Notas}
	\begin{minted}[bgcolor = DarkBackground, linenos=true]{python}
   notas = [7.5, 8.0, 6.5, 9.8, 7.0]

   soma = 0
   media = 0

   for nota in notas:
      somas+= nota

   print('Soma:', soma) # O '' entre a palavra 
   # indica que deve aparecer escrito
   
   media = soma/len(notas)

   print('Média: ', media)

   ## Alternativamente, o python tem algumas 
   # funções prontas para listas, por exemplo
   print('Soma: ', sum(notas))
   print('Maior nota: ', max(notas))
   print('Menor nota: ', min(notas))
 \end{minted}
\end{listing}

Se quisermos adicionar uma nota nova ao final da lista, utiliza-se o \texttt{append()}
\begin{minted}[bgcolor = DarkBackground, linenos=true]{python}
   notas = [7.5, 8.0, 6.5, 9.8, 7.0]

   notas.append(3.0)
   print(notas)
\end{minted}

É possível fazer uma ``lista de listas'', conhecida como matriz no python, que serão utilizadas na aula seguinte; por exemplo:
\begin{minted}[bgcolor = DarkBackground, linenos=true]{Python}
matriz = [[1,3,2,4], [2,4,6,8], [1,3,5,7,9]]

# Os elementos da matriz são acessados parecido com os da lista: 
# utilizando matriz[x][y], acessamos o elemento de índcie y na 
# x-ésima lista dentro da matriz. 

print(matriz[0][2]) # 3
matriz[0][2] = 9549 # Altera o segundo elemento da lista 0.

print(matriz[0][2]) # 9549
\end{minted}

\end{document}
